\scanpage{183}\mspage{38}%
% archive sheet number: 183
% Main writing: his hand. The archive pencil number is a separate hand.
% One continuous prose paragraph; L9 has no additional vertical gap or indentation.
dans cette première approche pré-catégorie dérivée, par la connaissance du groupe fondamental $\pi_1$, donc des catégories des coeff. \unsure{constant|constants} tordus $=\pi_1$-\unsure{module|modules}, les foncteurs $H^i(-,M)$ sur ceux-ci avec \unsure{la|le} structure de \unsure{foncteur|foncteurs} \deleted{\unsure{ab|ab.}, de \unsure{foncteur|foncteurs}} \inserted{\unsure{cohomologique|cohomologiques}} relatif aux suites exactes, la
% L7: an overwritten first layer is physically superposed on the beginning of structure; its letters and chronology remain uncertain.
\deleted{\unsure{struc}}structure de cup-produit etc --%
% L8: compact detached point at approximately (247,794), below the left end of the short dash; origin unassigned.
\raisebox{0pt}[0pt][0pt]{\parbox{0pt}{\figure{\begin{tikzpicture}[overlay]\fill (-0.8em,-0.55em) ellipse[x radius=0.30pt,y radius=0.40pt];\end{tikzpicture}}}}
reliés par certaines propriétés formelles. Une fois qu'on dispose du langage des cat. dérivées, la \inserted{sous catégorie de la catégorie dérivée des} complexes abéliens, \unsure{form-|formée|formés} \inserted{aux complexes de} faisceaux de coh. loc constante sur $X$, avec sa structure triangulé et sa structure \deleted{\unsure{sou|m|de}}multiplicative $\overset{\mathbb{L}}{\otimes}%
% L13: detached point at approximately (1165,1374), well below the tensor, is not an operator or index.
\raisebox{0pt}[0pt][0pt]{\parbox{0pt}{\figure{\begin{tikzpicture}[overlay]\fill (-0.35em,-0.9em) ellipse[x radius=0.32pt,y radius=0.48pt];\end{tikzpicture}}}}$
(et éventuellement $\mathbb{R}\mathrm{Hom}$ \ldots) donnant un candidat plus satisfaisant pour espérer récupérer le type d'homotopie. Je n'ai guère d'idée si%
% L17: small separate bent vertical stroke at approximately (241,1836), below si; origin unassigned.
\raisebox{0pt}[0pt][0pt]{\parbox{0pt}{\figure{\begin{tikzpicture}[overlay]\draw[line width=0.55pt] (-0.35em,-0.65em) -- (-0.29em,-0.80em) -- (-0.31em,-0.94em) -- (-0.24em,-0.97em);\end{tikzpicture}}}}
oui ou non ils%
% L17: small dark point with a coloured halo at approximately (802,1787), above-right of the final s; no textual role assigned.
\raisebox{0pt}[0pt][0pt]{\parbox{0pt}{\figure{\begin{tikzpicture}[overlay]\fill (0.15em,0.40em) circle[radius=0.30pt];\end{tikzpicture}}}}
% L17: separate irregular dark point at approximately (873,1810), immediately left of the p of permettent; origin unassigned.
\raisebox{0pt}[0pt][0pt]{\parbox{0pt}{\figure{\begin{tikzpicture}[overlay]\fill (0.50em,-0.18em) ellipse[x radius=0.35pt,y radius=0.45pt];\end{tikzpicture}}}}
permettent de le récupérer \inserted{en effet}, mais par contre%
% L18: two distinct points, approximately (944,1958) and (940,1965), below-left of contre; not a textual colon.
\raisebox{0pt}[0pt][0pt]{\parbox{0pt}{\figure{\begin{tikzpicture}[overlay]\fill (-2.1em,-0.70em) circle[radius=0.25pt];\fill (-2.2em,-0.90em) circle[radius=0.25pt];\end{tikzpicture}}}}
n'ai guère de doute qu'en allant au bout de la ``linéarisation'', c'est à dire en partant dans le
% L21: two physical underline segments: approximately (318,2273)--(558,2261), then (583,2262)--(672,2263). The second starts under the far-right return of l; the gap is preserved at letter granularity.
\uline{non abé}l\uline{ien} et travaillant avec la $(n+1)$-catégorie (\unsure{sans|sa|san} structures supplémentaires dessus!) des $n$-champs loc. constants
% L23: raised, heavily retraced sur X with a curved guide attaching after constants and before pour. X is clear; sur remains doubtful. No independent cancelling stroke is secure.
\inserted{\unsure{sur} $X$} pour tout $n$
